% Generated draft: requires mathematical and source review.
\documentclass[11pt]{article}
\usepackage[margin=25mm]{geometry}
\usepackage{fontspec}
\setmainfont{Latin Modern Roman}
\usepackage{amsmath,amssymb,mathtools,unicode-math}
\setmathfont{Latin Modern Math}
\usepackage{enumitem,needspace,xurl,hyperref}
\setlength{\parindent}{0pt}
\setlength{\parskip}{5pt}
\setlength{\emergencystretch}{3em}
\newcommand{\N}{\mathbb{N}}
\newcommand{\Z}{\mathbb{Z}}
\newcommand{\Q}{\mathbb{Q}}
\newcommand{\R}{\mathbb{R}}
\newcommand{\C}{\mathbb{C}}
\newcommand{\F}{\mathbb{F}}
\newcommand{\A}{\mathbb{A}}
\newcommand{\PP}{\mathbb{P}}
\newcommand{\E}{\mathbb{E}}
\newcommand{\eps}{\varepsilon}
\newcommand{\dd}{\,\mathrm{d}}
\begin{document}
\section*{UnsolvedMath ID: 30004833}
\textbf{Source problem code:} catalogue entry
\subsection{Definitions and mathematical statement}
Write \(n(B)=\dim B\) and \(M(B)=L^\infty(G^+(B))\), using the source notation. Under its conventional interpretation, \(G^+(B)\) is a quantum automorphism group associated with \(B\), and \(L^\infty(G^+(B))\) is the von Neumann algebra obtained from its Haar-state GNS representation. However, the supplied excerpt does not define the admissible algebras \(B\), the ground field for their dimension, or any state, trace, or normalization required to define \(G^+(B)\). Thus these conventions cannot be treated as recovered source hypotheses.

Let \(\mathcal B\) stand provisionally for the intended class of admissible objects, including any required auxiliary data. With \(\cong\) interpreted as a normal unital \(*\)-isomorphism of von Neumann algebras, the intended dimension-invariance question is whether
\[
\forall n\in\N,\quad n\geq5\ \Longrightarrow\ 
\forall B_1,B_2\in\mathcal B,\quad
\bigl(n(B_1)=n(B_2)=n\bigr)
\ \Longrightarrow\ M(B_1)\cong M(B_2).
\]
This is a conditional formulation: the class \(\mathcal B\), its auxiliary data, and the intended isomorphism convention require confirmation from the report. No trace-preserving condition or factoriality assertion is added.

The original excerpt also asks whether its Theorem C can be used to establish this dimension invariance. Its surrounding discussion concerns commuting group actions and isomorphisms of iterated algebraic crossed products involving \(\mathcal O(G^+(B))\), \(\mathcal O(S_n^+)\), and matrix-algebra tensor factors. The supplied transcription is corrupted and refers to notation and hypotheses from an absent Theorem A; it does not provide a reliable complete statement of Theorem C. The suggested use of that theorem is retained as source context, without asserting that these algebraic isomorphisms imply the displayed von Neumann algebra isomorphism.
\subsection{Short English statement}
For each integer \(n\geq5\), are the von Neumann algebras \(L^\infty(G^+(B))\) isomorphic for every admissible \(B\) of dimension \(n\)? The source suggests using Theorem C. The admissibility hypotheses and auxiliary data defining \(G^+(B)\) are missing, so this remains a conditional statement requiring human mathematical review.
\subsection{Sources}
\begin{itemize}
\item Local registry: supplied UnsolvedMath record, permanent ID \(30004833\), exact source problem code \texttt{catalogue entry}; project registry locations are \texttt{data/unsolvedmath/problems.json} and \texttt{catalog/problems.json}. This draft uses the supplied complete record without inspecting those files.
\item Supplied report citation: \emph{Quantum Groups -- Algebra, Analysis and Category Theory} (2022), Oberwolfach Reports, pp.\ 2397--2458, DOI \texttt{10.4171/owr/2021/44}. Supplied source URL: \url{https://doi.org/10.4171/owr/2021/44}. No individual proposer is identified in the record; its proposed year is 2021.
\item Supplied literature reference: ``Oberwolfach Report 44,'' \url{https://doi.org/10.4171/owr/2021/44}. The literature assessment checked at \texttt{2026-08-22} labels the problem open and reports that the curation search found no primary result establishing dependence solely on \(\dim B\). This is supplied dated curation, not an independently verified literature assessment.
\item The registry labels the statement \texttt{corrected\_verified} and says that the linked report was inspected to isolate the explicit question from the surrounding theorem and crossed-product formulas. These are supplied verification claims; this draft does not independently verify them or claim to have visited the report.
\end{itemize}
\end{document}
